\documentclass[11pt,a4paper]{article}

%% ── Packages ──────────────────────────────────────────────────────────────
\usepackage[utf8]{inputenc}
\usepackage[T1]{fontenc}
\usepackage[english]{babel}
\usepackage{geometry}
\geometry{margin=2.5cm}
\usepackage{parskip}
\usepackage{enumitem}
\usepackage{hyperref}
\hypersetup{colorlinks=true, linkcolor=blue!60!black, urlcolor=blue!60!black, citecolor=blue!60!black}
\usepackage{booktabs}
\usepackage{graphicx}
\usepackage{amsmath,amssymb}
\usepackage{natbib}
\usepackage{xcolor}
\usepackage{fancyhdr}
\usepackage{titlesec}
\usepackage{tabularx}

%% ── Page style ────────────────────────────────────────────────────────────
\pagestyle{fancy}
\fancyhf{}
\lhead{\small Quantitative Methods in Finance --- Final Exam 2025--2026}
\rhead{\small Master 2 Research}
\cfoot{\thepage}

\titleformat{\section}{\large\bfseries}{\thesection.}{0.5em}{}
\titleformat{\subsection}{\normalsize\bfseries}{\thesubsection.}{0.5em}{}

%% ── Title ─────────────────────────────────────────────────────────────────
\title{%
  \vspace{-1.5cm}
  {\large Quantitative Methods in Finance}\\[0.3em]
  {\Large\bfseries Final Exam --- Take-Home}\\[0.3em]
  {\large\itshape Counterfactual Inflation Analysis:\\
   What If Ukraine Had Been Part of the Euro Area?}
}
\author{%
  Eric Vansteenberghe\\[0.2em]
  \small Banque de France \& Universit\'{e} Paris 1 Panth\'{e}on-Sorbonne
}
\date{Academic Year 2025--2026}

\begin{document}
\maketitle
\thispagestyle{fancy}

%% ══════════════════════════════════════════════════════════════════════════
\section*{Context and Motivation}
%% ══════════════════════════════════════════════════════════════════════════

On 23 June 2022, the European Council granted Ukraine official candidate status
for accession to the European Union. While EU membership and Euro Area
membership are distinct steps---separated by years of institutional convergence
and the fulfilment of the Maastricht criteria---the question of whether Ukraine
could eventually adopt the euro is no longer purely hypothetical. It is a
question that central banks, including the ECB and the national central banks of
the Eurosystem, will increasingly need to assess.

This exam asks you to engage with a foundational question at the intersection of
monetary economics and applied econometrics:

\begin{center}
\itshape
What would Ukraine's inflation trajectory have looked like\\
had it been a member of the Euro Area?
\end{center}

This is a \textbf{counterfactual} question. Ukraine was never in the Euro Area;
it maintained its own currency (the hryvnia), its own central bank (the National
Bank of Ukraine, NBU), and its own monetary policy throughout the sample period.
Your task is to construct a plausible alternative inflation path for Ukraine
under the hypothetical scenario of Euro Area membership, using the econometric
methods studied in this course.

%% ══════════════════════════════════════════════════════════════════════════
\section*{Why This Question Matters}
%% ══════════════════════════════════════════════════════════════════════════

The academic literature on \textbf{Optimal Currency Areas} (OCA), initiated by
\citet{mundell1961} and extended by \citet{mckinnon1963} and \citet{kenen1969},
identifies the fundamental trade-off faced by a country considering membership in
a monetary union: the \emph{loss of independent monetary and exchange-rate
policy} versus the \emph{gains from reduced transaction costs, price
transparency, and imported monetary credibility}.

For Ukraine, this trade-off is particularly stark. Consider the following
channels through which Euro Area membership would have reshaped the inflation
process:

\begin{enumerate}[leftmargin=2em]

  \item \textbf{Loss of the exchange-rate buffer.}
  Ukraine experienced three major hryvnia devaluations: approximately $-50\%$
  against the US dollar in 2008--2009 (global financial crisis), over $-200\%$
  in 2014--2015 (Crimea annexation and Donbas conflict), and approximately
  $-25\%$ in 2022 (full-scale invasion). These devaluations were a critical
  adjustment mechanism. Under Euro Area membership, this channel would have been
  \emph{unavailable}. The real adjustment would have had to occur through
  internal devaluation---wage and price deflation---as experienced by Greece,
  Ireland, and the Baltic states during the 2010--2012 sovereign debt crisis.

  \item \textbf{Imported monetary credibility.}
  The ECB's inflation-targeting framework and institutional independence provide
  a nominal anchor that the NBU has historically lacked. Before the adoption of
  inflation targeting in 2016, Ukraine operated under various exchange-rate and
  monetary-aggregate regimes with limited credibility, leading to persistent
  inflation expectations well above Euro Area levels. Euro Area membership would
  have \emph{imported} the ECB's credibility, likely anchoring inflation
  expectations at lower levels---the so-called \citet{frankel_rose1998}
  endogeneity hypothesis, whereby joining a currency union makes a country
  \emph{more} suitable for it ex post.

  \item \textbf{One monetary policy for heterogeneous shocks.}
  The ECB sets a single policy rate for the entire Euro Area. Ukraine's economy
  is structurally different from the Euro Area core: it is more exposed to
  energy-price shocks (historically dependent on Russian natural gas), more
  reliant on agricultural commodity exports, and subject to geopolitical shocks
  that are largely asymmetric---i.e., they do not affect Euro Area members
  similarly. A single monetary policy calibrated for the Euro Area aggregate
  might not have been \textit{optimal} for Ukraine's specific
  conditions.

  \item \textbf{Fiscal discipline and the Stability and Growth Pact.}
  Euro Area membership comes with fiscal constraints (the 3\% deficit and 60\%
  debt-to-GDP rules, however imperfectly enforced). These constraints would have
  limited Ukraine's ability to use fiscal policy as a substitute for the lost
  monetary-policy and exchange-rate instruments---amplifying the cost of
  asymmetric shocks but potentially reducing the inflation bias associated with
  deficit monetisation.

  \item \textbf{Financial integration and capital flows.}
  Euro Area membership would have eliminated currency risk for foreign investors,
  likely resulting in lower sovereign borrowing costs and larger capital inflows
  \emph{in normal times}---but also removing the ability to impose capital
  controls or let the exchange rate depreciate during sudden stops, as Ukraine
  did in 2008--2009, 2014--2015, and 2022.

\end{enumerate}

The counterfactual you construct is therefore not merely a statistical exercise.
It embodies a set of \textbf{structural hypotheses} about how the inflation
process would have differed under an alternative monetary regime. The quality of
your answer depends on how explicitly and coherently you operationalise these
hypotheses in your econometric framework.


%% ══════════════════════════════════════════════════════════════════════════
\section*{Part A --- Ukraine's Monetary Regime: A Preliminary Analysis (30\%)}
%% ══════════════════════════════════════════════════════════════════════════

Before constructing a counterfactual, you must understand \emph{what Ukraine
actually did} with its monetary sovereignty. A na\"{\i}ve framing assumes
Ukraine operated under a freely floating exchange rate with full monetary policy
independence throughout the sample period. This is incorrect.

\subsection*{Task}

Construct a \textbf{documented chronology} of the National Bank of Ukraine's
exchange rate regime from 2000 to 2025. Your chronology must identify:

\begin{enumerate}[label=(\alph*), leftmargin=2em]

  \item The periods during which the NBU maintained a \textbf{de facto peg} to
  the US dollar (and the approximate level of the peg);

  \item The \textbf{major devaluation episodes}, with approximate dates,
  magnitudes, and triggers;

  \item The periods during which significant \textbf{capital controls} were in
  place;

  \item The transition to \textbf{inflation targeting} (official date, target
  path, and whether the NBU's behaviour actually changed in practice);

  \item For each period, your assessment of the \textbf{IMF's de facto
  classification} of Ukraine's exchange rate arrangement (stabilised
  arrangement, managed float, conventional peg, etc.).

\end{enumerate}

\subsection*{Deliverables for Part~A}

\begin{enumerate}[label=\textbf{\arabic*.}, leftmargin=2em]

  \item A \textbf{summary table} with columns: period, approximate UAH/USD
  rate, de facto regime classification, and key events.

  \item A \textbf{brief argument} (one paragraph, in your code or a separate
  text file) answering: \emph{During which periods did Ukraine have genuine
  monetary sovereignty? During which periods was its monetary policy de facto
  constrained by the exchange-rate peg? What does this imply for the
  ``treatment'' that Euro Area membership would represent?}

\end{enumerate}

\subsection*{Why this matters}

This preliminary analysis is not optional background---it is central to your
identification strategy in Part~B. If Ukraine was already operating a de facto
dollar peg for most of the sample, then Euro Area membership would not have
represented a large change in monetary regime during the peg periods. The
``treatment'' of joining the Euro Area is concentrated in the episodes where
Ukraine \emph{deviated} from its peg---the devaluations and the post-2016
inflation-targeting period. Your counterfactual must reflect this asymmetry.

\paragraph{Sources.} You should consult IMF as well as NBU publications and academic sources. Cite your sources.


%% ══════════════════════════════════════════════════════════════════════════
\section*{Data}
%% ══════════════════════════════════════════════════════════════════════════

\subsection*{Repository data}

The following datasets are available in the GitHub repository:

\begin{center}
\url{https://github.com/skimeur/pioneer-detection-method}
\end{center}

\begin{itemize}[leftmargin=2em]
  \item \texttt{data\_ecb\_hicp\_panel.csv} --- Monthly Harmonised Index of
    Consumer Prices (HICP), year-on-year percentage change, for 11 Euro Area
    countries (Austria, Belgium, Germany, Spain, Finland, France, Greece,
    Ireland, Italy, Netherlands, Portugal), from January 2000 to December 2025.
    Source: ECB Data Portal.
    \begin{itemize}
    \item Nota bene: you might need to augment this list of countries if you do a synthetic control.
    \end{itemize}

  \item \texttt{data\_ukraine\_cpi\_raw.csv} --- Monthly Consumer Price Index
    for Ukraine, expressed as an \textbf{index with the previous month as base}
    (i.e., a value of 101.5 means prices rose 1.5\% month-on-month), from
    February 2005 to December 2025. Source: State Statistics Service of Ukraine
    (SSSU) via SDMX.
\end{itemize}

\paragraph{Important note on data harmonisation.}
The two datasets are \emph{not} in the same format. The Euro Area data are
year-on-year inflation rates (\%), while the Ukrainian data are month-on-month
price indices. You will need to transform the Ukrainian series into year-on-year
inflation rates to make it comparable with the Euro Area panel.

\subsection*{External data: a mandatory requirement}

\textbf{You are required to augment the repository data with additional
macroeconomic time series from external sources.} The methods studied in this
course---structural VARs, local projections---require variables beyond inflation
alone. At a minimum, you should consider sourcing:

\begin{itemize}[leftmargin=2em]
  \item \textbf{Real GDP or industrial production} for Ukraine and Euro Area
    countries (quarterly or monthly), from Eurostat, the IMF's International
    Financial Statistics (IFS), or the World Bank;
  \item \textbf{Exchange rate series} (UAH/USD, UAH/EUR) from the NBU or IMF;
  \item Any other macroeconomic variables your identification strategy requires
    (policy rates, commodity prices, fiscal balances, etc.).
\end{itemize}

The quality, relevance, and documentation of your external data choices are
\textbf{part of the grading}. You must cite your data sources and explain why
each variable is needed for your econometric framework.


%% ══════════════════════════════════════════════════════════════════════════
\section*{Key Literature for the Counterfactual}
%% ══════════════════════════════════════════════════════════════════════════

The following papers are selected not as a general reading list, but as
\textbf{direct methodological inputs} to the construction of your
counterfactual. Each one offers a specific tool, identification strategy, or
conceptual framework that you can adapt. You are not required to read all of
them---choose those relevant to your chosen approach.


%% ── 1. Shock decomposition ───────────────────────────────────────────────
\subsection*{Structural shock decomposition: Bayoumi and Eichengreen (1993)}

\citet{bayoumi_eichengreen1993} estimate bivariate structural VARs (output
growth, inflation) for pairs of countries, using the \citet{blanchard_quah1989}
long-run identification: supply shocks have permanent effects on output, demand
shocks do not. They then compare the \emph{correlation of structural shocks}
across countries to assess OCA suitability, distinguishing a European ``core''
(whose shocks are symmetric) from a ``periphery'' (whose shocks are
idiosyncratic).

\paragraph{How to use it for the counterfactual.}
You can replicate the Bayoumi--Eichengreen decomposition for Ukraine and each
Euro Area country, then construct the counterfactual via \emph{shock
replacement}: if Ukraine had been in the Euro Area, its \textbf{demand shocks}
would have been determined by ECB monetary policy (i.e., the Euro Area aggregate
demand shocks), while its \textbf{supply shocks}---driven by real factors
(energy prices, agriculture, geopolitics)---would have remained its own. The
counterfactual inflation path is obtained by feeding Ukraine's supply shocks and
the Euro Area's demand shocks through Ukraine's (or the Euro Area's) estimated
impulse response functions. This requires both output and inflation data, which
is why external data sourcing is mandatory.


%% ── 2. Common factor ─────────────────────────────────────────────────────
\subsection*{Common inflation factors: Ciccarelli and Mojon (2010)}

\citet{ciccarelli_mojon2010} show that a large fraction of inflation variation
across countries can be explained by a single common factor, which they
interpret as ``global inflation.'' Their approach extracts a common component
from a panel of inflation series using dynamic factor model techniques
(principal components or state-space estimation).

\paragraph{How to use it for the counterfactual.}
Estimate the Euro Area common inflation factor from the panel of 11 countries.
Then model Ukraine's inflation as:
\[
  \pi_{UKR,t} \;=\; \lambda_{UKR}\, F_t^{EA} \;+\; \varepsilon_{UKR,t}
\]
where $F_t^{EA}$ is the Euro Area common factor, $\lambda_{UKR}$ is Ukraine's
loading on that factor (estimated from ``quiet'' periods when Ukraine's
inflation co-moved with Europe), and $\varepsilon_{UKR,t}$ is the
Ukraine-specific residual. The counterfactual ``Ukraine-in-the-Euro-Area''
inflation is $\hat{\pi}_{UKR,t}^{CF} = \hat{\lambda}_{UKR}\, \hat{F}_t^{EA}$:
only the common European component, without the idiosyncratic shocks that
Ukraine's own monetary policy and exchange rate produced. This is a panel
method that can be implemented with the inflation data alone; adding
macroeconomic covariates strengthens identification.


%% ── 3. Synthetic control ─────────────────────────────────────────────────
\subsection*{Synthetic control: Abadie, Diamond, and Hainmueller (2010)}

\citet{abadie_diamond_hainmueller2010} propose constructing a ``synthetic''
version of a treated unit as a convex combination of untreated donor units,
with weights chosen to match the treated unit's pre-treatment trajectory.

\paragraph{A difficulty for this exercise.}
In the classic synthetic control setup, there is a clean treatment date (before
which the treated unit and its synthetic counterpart should track each other,
and after which they diverge). For Ukraine, \textbf{there is no clean treatment
date}: Ukraine never joined the Euro Area, so there is no before/after. You may
attempt to invert the logic---treating Ukraine's \emph{absence} from the Euro
Area as the treatment---but this requires careful conceptual justification.

\paragraph{A note on recent Euro Area accessions.}
Since 1999, ten countries have adopted the euro. Several of them---particularly
among the post-Soviet Baltic states---share structural features with Ukraine
(small, open, post-Soviet economies with a history of exchange-rate pegs).
Their accession dates and pre/post inflation paths may contain useful
information for your analysis:

\begin{center}
\small
\begin{tabular}{llll}
\toprule
\textbf{Country} & \textbf{EU accession} & \textbf{Euro adoption} & \textbf{Notes} \\
\midrule
Greece    & 1981  & 1 Jan 2001 & Debt crisis 2010--2015 \\
Slovenia  & 2004  & 1 Jan 2007 & First 2004-enlargement country \\
Cyprus    & 2004  & 1 Jan 2008 & Banking crisis 2012--2013 \\
Malta     & 2004  & 1 Jan 2008 & Small island economy \\
Slovakia  & 2004  & 1 Jan 2009 & Adopted euro during GFC \\
Estonia   & 2004  & 1 Jan 2011 & Post-Soviet; currency board pre-euro \\
Latvia    & 2004  & 1 Jan 2014 & Post-Soviet; severe 2008--09 recession \\
Lithuania & 2004  & 1 Jan 2015 & Post-Soviet; last Baltic to join \\
Croatia   & 2013  & 1 Jan 2023 & De facto euro-linked kuna pre-adoption \\
Bulgaria  & 2007  & 1 Jan 2026 & Currency board since 1997 \\
\bottomrule
\end{tabular}
\end{center}


%% ── 4. Fear of floating ──────────────────────────────────────────────────
\subsection*{Constrained sovereignty: Calvo and Reinhart (2002)}

\citet{calvo_reinhart2002} document that many countries that officially declare
a floating exchange rate in practice intervene heavily to limit exchange-rate
movements---a pattern they call ``fear of floating.'' The implication is that
\emph{de jure} monetary sovereignty can be very different from \emph{de facto}
monetary sovereignty.

\paragraph{Why this is critical for Ukraine.}
As your Part~A analysis should reveal, the NBU maintained a de facto peg to the
US dollar for most of the period 2000--2014, and returned to a fixed rate under
wartime conditions in 2022--2023. During peg periods, Ukraine had \emph{no
independent monetary policy} in practice---the ``impossible trinity'' was
binding. This means that the counterfactual ``treatment'' of Euro Area
membership is \textbf{not constant over time}: it is large during episodes when
Ukraine exercised genuine monetary autonomy (the devaluations, the post-2016
inflation-targeting period) and small during peg periods (when monetary policy
was already subordinated to the exchange-rate anchor). Your identification
strategy should account for this time-varying treatment intensity.


%% ── 5. Credibility ──────────────────────────────────────────────────────
\subsection*{Credibility import: Giavazzi and Pagano (1988); Barro and Gordon (1983)}

\citet{giavazzi_pagano1988} argue that countries with weak monetary credibility
can ``tie their hands'' by joining a fixed exchange-rate arrangement (or a
monetary union), importing the credibility of the anchor country's central
bank. \citet{barro_gordon1983} provide the theoretical foundation: under
discretion, a central bank faces a time-consistency problem that generates an
inflation bias; an external commitment device eliminates this bias.

\paragraph{Relevance for Ukraine.}
The NBU's credibility was low before 2016 (no formal inflation target, repeated
peg collapses). After adopting inflation targeting in 2016, credibility improved
but remained fragile. A \textbf{sanity check} for your counterfactual: the gap
between actual Ukrainian inflation and your counterfactual Euro Area path
should be \emph{larger} in the pre-2016 period (when credibility gains from Euro
Area membership would have been substantial) and potentially \emph{smaller}
after 2016 (when the NBU had partially converged toward a modern
inflation-targeting framework). If your counterfactual does not exhibit this
pattern, reconsider your assumptions.


%% ── 6. Sovereign fragility ──────────────────────────────────────────────
\subsection*{Sovereign fragility in monetary unions: De~Grauwe (2012)}

\citet{degrauwe2012} argues that members of a monetary union are uniquely
vulnerable to self-fulfilling sovereign debt crises, because they cannot rely
on their own central bank as a lender of last resort in the currency of
denomination. This fragility was illustrated by the sovereign spreads divergence
in 2010--2012 and the eventual need for the ECB's Outright Monetary
Transactions (OMT) programme.

\paragraph{Relevance for Ukraine.}
Had Ukraine been in the Euro Area during the 2014 or 2022 crises, it would have
avoided a \emph{currency crisis} but might have faced a \emph{sovereign debt
crisis} instead---with no national central bank to provide emergency liquidity.
The counterfactual is therefore not simply ``lower inflation'' but potentially
``a different type of macroeconomic crisis.'' This is not a required modelling
element, but acknowledging this channel in your interpretation will demonstrate
deeper understanding of the trade-offs.


%% ══════════════════════════════════════════════════════════════════════════
\section*{Part B --- The Counterfactual (50\%)}
%% ══════════════════════════════════════════════════════════════════════════

Using the methods studied in this course---and any additional econometric tools
you deem appropriate---construct a counterfactual inflation time series for
Ukraine under the hypothesis that it had been a member of the Euro Area.

\subsection*{Deliverables for Part~B}

\begin{enumerate}[label=\textbf{\arabic*.}, leftmargin=2em]

  \item \textbf{A single figure} showing two time series on the same axes:
  \begin{itemize}
    \item Ukraine's \emph{actual} year-on-year inflation (computed from the raw data);
    \item Your \emph{counterfactual} ``Ukraine-in-the-Euro-Area'' inflation.
  \end{itemize}
  The figure should be clearly labelled, with a legend, axis titles, and
  appropriate date formatting.

  \item \textbf{A brief interpretation} (one paragraph) answering: \emph{What
  does your counterfactual imply about the cost or benefit of monetary
  sovereignty for Ukraine, particularly during the 2008--2009, 2014--2015, and
  2022 crises?}

  \item \textbf{Your code}, which must be fully reproducible. All external data
  must either be downloaded programmatically within your code or provided
  alongside it with clear documentation.

\end{enumerate}

\subsection*{Methodological requirements}

\begin{itemize}[leftmargin=2em]

  \item Your method must reflect an \textbf{explicit identification strategy}:
  you must articulate (through your code structure and comments) what ``being in
  the Euro Area'' means in terms of your model's parameters, shocks, or
  constraints.

  \item \textbf{Your counterfactual must differ from the simple cross-sectional
  mean of Euro Area countries' inflation.} This constraint ensures that your
  approach involves genuine econometric reasoning rather than a mechanical
  average.

  \item Your identification strategy must be \textbf{consistent with your
  Part~A analysis}. If you found that Ukraine had no genuine monetary
  sovereignty during peg periods, your counterfactual should reflect a smaller
  ``treatment effect'' during those periods.

  \item You are expected to use \textbf{structural VARs, local projections, or
  comparable methods} studied in this course. You may complement them with panel
  methods (factor models, synthetic control), but the core of your analysis
  should demonstrate mastery of the SVAR/LP toolkit.

\end{itemize}


%% ══════════════════════════════════════════════════════════════════════════
\section*{Grading Criteria}
%% ══════════════════════════════════════════════════════════════════════════

\begin{center}
\small
\begin{tabularx}{\textwidth}{lrX}
\toprule
\textbf{Component} & \textbf{Weight} & \textbf{Criteria} \\
\midrule
\textbf{Part A:} Monetary regime & 30\% &
  Accuracy of the chronology; quality of sources; coherence of the argument
  about which periods represent genuine monetary sovereignty; correct
  identification of de facto vs.\ de jure regimes. \\[0.8em]
\textbf{Part B:} Identification & 30\% &
  Is the counterfactual grounded in a coherent economic and econometric
  framework? Are the assumptions explicit? Does the method go beyond mechanical
  projection? Is it consistent with Part~A? \\[0.8em]
\textbf{Part B:} Technical execution & 20\% &
  Is the code correct, reproducible, and well-structured? Is the data
  transformation handled properly? Are estimation choices (lag selection,
  stationarity testing, external data sourcing) justified? \\[0.8em]
\textbf{Part B:} Output and interpretation & 20\% &
  Is the figure informative and professionally presented? Does the
  counterfactual exhibit plausible dynamics? Does the interpretation demonstrate
  economic understanding of the OCA trade-offs? \\
\bottomrule
\end{tabularx}
\end{center}


%% ══════════════════════════════════════════════════════════════════════════
\section*{Practical Information}
%% ══════════════════════════════════════════════════════════════════════════

\begin{itemize}[leftmargin=2em]
  \item \textbf{Format:} take-home exam, individual work.
  \item \textbf{Duration:} you have until \textbf{1st of May 2026} to submit.
  \item \textbf{Submission:} upload your code, output figure(s), and Part~A
    deliverables to \textbf{GitHub}.
  \item \textbf{Language:} any programming language (Python, R, Julia,
    MATLAB, \ldots). The code must run from the repository data and your
    documented external data without manual intervention.
  \item \textbf{External data:} all external datasets must be either
    (a)~downloaded programmatically within your code, or (b)~provided as files
    alongside your submission with full source documentation.
  \item \textbf{Collaboration policy:} this is individual work. You may
    consult any published reference, documentation, or AI tool, but the
    econometric reasoning and code must be your own. Identical or
    near-identical submissions will be treated as plagiarism.
\end{itemize}

%% ══════════════════════════════════════════════════════════════════════════
\section*{Suggested References}
%% ══════════════════════════════════════════════════════════════════════════

The following references supplement those discussed in the literature section
above. Choose those relevant to your chosen method.

\begin{description}[leftmargin=1em, style=nextline]

  \item[Optimal Currency Areas]
  \citet{mundell1961}; \citet{mckinnon1963}; \citet{kenen1969};
  \citet{frankel_rose1998}.

  \item[Structural VARs and local projections]
  \citet{sims1980}; \citet{blanchard_quah1989}; \citet{stock_watson2001};
  \citet{jorda2005}; \citet{kilian_lutkepohl2017}.

  \item[Counterfactual methods]
  \citet{abadie_diamond_hainmueller2010} (synthetic control);
  \citet{angrist_pischke2009} (DiD and potential outcomes).

  \item[Euro Area inflation dynamics]
  \citet{ciccarelli_mojon2010}.

  \item[Monetary credibility and exchange-rate regimes]
  \citet{barro_gordon1983}; \citet{giavazzi_pagano1988};
  \citet{calvo_reinhart2002}; \citet{degrauwe2012}.

\end{description}


%% ══════════════════════════════════════════════════════════════════════════
%% Bibliography
%% ══════════════════════════════════════════════════════════════════════════

\begingroup
\renewcommand{\section}[2]{}  % suppress "References" heading from natbib
\begin{thebibliography}{99}
\small

\bibitem[Abadie et~al.(2010)]{abadie_diamond_hainmueller2010}
Abadie, A., Diamond, A., and Hainmueller, J. (2010).
\newblock Synthetic control methods for comparative case studies: Estimating the
  effect of California's tobacco control program.
\newblock \emph{Journal of the American Statistical Association},
  105(490):493--505.

\bibitem[Angrist and Pischke(2009)]{angrist_pischke2009}
Angrist, J.~D. and Pischke, J.-S. (2009).
\newblock \emph{Mostly Harmless Econometrics: An Empiricist's Companion}.
\newblock Princeton University Press.

\bibitem[Barro and Gordon(1983)]{barro_gordon1983}
Barro, R.~J. and Gordon, D.~B. (1983).
\newblock Rules, discretion and reputation in a model of monetary policy.
\newblock \emph{Journal of Monetary Economics}, 12(1):101--121.

\bibitem[Bayoumi and Eichengreen(1993)]{bayoumi_eichengreen1993}
Bayoumi, T. and Eichengreen, B. (1993).
\newblock Shocking aspects of European monetary integration.
\newblock In Torres, F. and Giavazzi, F., editors, \emph{Adjustment and Growth
  in the European Monetary Union}, pages 193--229. Cambridge University Press.

\bibitem[Blanchard and Quah(1989)]{blanchard_quah1989}
Blanchard, O.~J. and Quah, D. (1989).
\newblock The dynamic effects of aggregate demand and supply disturbances.
\newblock \emph{American Economic Review}, 79(4):655--673.

\bibitem[Calvo and Reinhart(2002)]{calvo_reinhart2002}
Calvo, G.~A. and Reinhart, C.~M. (2002).
\newblock Fear of floating.
\newblock \emph{The Quarterly Journal of Economics}, 117(2):379--408.

\bibitem[Ciccarelli and Mojon(2010)]{ciccarelli_mojon2010}
Ciccarelli, M. and Mojon, B. (2010).
\newblock Global inflation.
\newblock \emph{The Review of Economics and Statistics}, 92(3):524--535.

\bibitem[De~Grauwe(2012)]{degrauwe2012}
De~Grauwe, P. (2012).
\newblock The governance of a fragile eurozone.
\newblock \emph{Australian Economic Review}, 45(3):255--268.

\bibitem[Frankel and Rose(1998)]{frankel_rose1998}
Frankel, J.~A. and Rose, A.~K. (1998).
\newblock The endogeneity of the optimum currency area criteria.
\newblock \emph{The Economic Journal}, 108(449):1009--1025.

\bibitem[Giavazzi and Pagano(1988)]{giavazzi_pagano1988}
Giavazzi, F. and Pagano, M. (1988).
\newblock The advantage of tying one's hands: EMS discipline and central bank
  credibility.
\newblock \emph{European Economic Review}, 32(5):1055--1075.

\bibitem[Jord\`{a}(2005)]{jorda2005}
Jord\`{a}, \`{O}. (2005).
\newblock Estimation and inference of impulse responses by local projections.
\newblock \emph{American Economic Review}, 95(1):161--182.

\bibitem[Kenen(1969)]{kenen1969}
Kenen, P.~B. (1969).
\newblock The theory of optimum currency areas: An eclectic view.
\newblock In Mundell, R.~A. and Swoboda, A.~K., editors, \emph{Monetary
  Problems of the International Economy}, pages 41--60. University of Chicago
  Press.

\bibitem[Kilian and L\"{u}tkepohl(2017)]{kilian_lutkepohl2017}
Kilian, L. and L\"{u}tkepohl, H. (2017).
\newblock \emph{Structural Vector Autoregressive Analysis}.
\newblock Cambridge University Press.

\bibitem[McKinnon(1963)]{mckinnon1963}
McKinnon, R.~I. (1963).
\newblock Optimum currency areas.
\newblock \emph{American Economic Review}, 53(4):717--725.

\bibitem[Mundell(1961)]{mundell1961}
Mundell, R.~A. (1961).
\newblock A theory of optimum currency areas.
\newblock \emph{American Economic Review}, 51(4):657--665.

\bibitem[Sims(1980)]{sims1980}
Sims, C.~A. (1980).
\newblock Macroeconomics and reality.
\newblock \emph{Econometrica}, 48(1):1--48.

\bibitem[Stock and Watson(2001)]{stock_watson2001}
Stock, J.~H. and Watson, M.~W. (2001).
\newblock Vector autoregressions.
\newblock \emph{Journal of Economic Perspectives}, 15(4):101--115.

\end{thebibliography}
\endgroup

\end{document}
